\chapter{P}
\label{chap:p}

\term{P-adic Number}
Elements of the field $\mathbb{Q}_p$, constructed as the completion of the rational numbers $\mathbb{Q}$ with respect to the $p$-adic valuation $|x|_p = p^{-v_p(x)}$. These numbers are fundamental in modern number theory and arithmetic geometry.

\term{P-adic Valuation}
A function $v_p: \mathbb{Q} \to \mathbb{Z} \cup \{\infty\}$ where $v_p(n)$ is the exponent of the highest power of $p$ that divides $n$. It satisfies the ultrametric inequality: $v_p(x+y) \geq \min(v_p(x), v_p(y))$.

\term{Painter's Paradox}
The observation connected with Gabriel's horn: the horn can be filled with a finite volume of paint, yet covering its surface seems to require an infinite amount. Under an idealized model with paint of no thickness the two statements conflict; in practice a finite film of paint on an infinite area still uses infinite volume, and the paradox highlights the difference between surface area and volume.

\term{Palindromic Sequence}
A sequence of symbols or numbers that reads the same forwards and backwards. In combinatorics on words, the study of palindromes relates to the structure of free monoids and Sturmian words.

\term{Parabolic PDE}
A partial differential equation of heat type, with one distinguished direction of time and smoothing evolution.

\term{Parallelepiped}
A three-dimensional figure whose six faces are all parallelograms. Its volume is given by the absolute value of the scalar triple product: $V = | \mathbf{a} \cdot (\mathbf{b} \times \mathbf{c}) |$.

\term{Parallel Transport}
A way of moving a vector along a curve on a manifold such that the vector "remains parallel" to itself, meaning its covariant derivative along the curve is zero. The result depends on the path taken if the manifold has curvature.

\term{Parameter}
A quantity that can be varied and appears as a constant in an equation, or the coordinate of a parametric curve.

\term{Parametric Equation}
An equation describing a curve through functions of one or more parameters, such as $x = \cos t$, $y = \sin t$.

\term{Parity}
The evenness or oddness of an integer; also the property of a function that is even or odd.

\term{Parrondo's Paradox}
Two games that are each losing can be combined, played alternately or by a random rule, into a winning strategy. It arises from capital-dependent or history-dependent strategies and connects to the theory of Brownian ratchets in physics; the winning arises from the nonlinearity of the combined process.

\term{Parseval's Identity}
The equality $\sum |a_n|^2 = \int |f|^2$ relating the Fourier coefficients of a function to its energy.

\term{Partial Derivative}
The derivative of a multivariable function with respect to one variable, while holding all other variables constant. It measures the rate of change along a specific coordinate axis.

\term{Partial Fraction}
The decomposition of a rational function into a sum of simpler rational terms, useful for integration.

\term{Partial Sum}
The sum of the first $n$ terms of a series, $S_n = a_1 + \cdots + a_n$.

\term{Partial Order}
A binary relation $\leq$ on a set $S$ that is reflexive ($a \leq a$), antisymmetric ($a \leq b \land b \leq a \implies a = b$), and transitive ($a \leq b \land b \leq c \implies a \leq c$).

\term{Partition (Integer)}
A way of writing a positive integer $n$ as a sum of positive integers, where the order of the summands does not matter. The partition function $p(n)$ grows rapidly and is studied using generating functions and Hardy-Ramanujan asymptotics.

\term{Path-Connected Space}
A topological space $X$ where any two points $x, y \in X$ can be joined by a continuous path $\gamma: [0,1] \to X$ such that $\gamma(0)=x$ and $\gamma(1)=y$. Path-connectedness implies connectedness.

\term{Path Integral}
A formulation of quantum mechanics developed by Richard Feynman, where the probability amplitude for a particle to travel from $a$ to $b$ is the sum over all possible paths, weighted by $e^{iS/\hbar}$ where $S$ is the action.

\term{Pattern}
A regular, repeating or predictable structure, studied formally in combinatorics and formal language theory.

\term{Pearson Correlation Coefficient}
A measure of the linear correlation between two variables, $\rho_{XY} = \operatorname{Cov}(X,Y) / (\sigma_X \sigma_Y)$, taking values in $[-1, 1]$. It is the sample correlation used in linear regression.

\term{Perfect Number}
A positive integer that is equal to the sum of its proper positive divisors. Examples include 6, 28, and 496. All known perfect numbers are even and given by the formula $2^{p-1}(2^p - 1)$ where $2^p-1$ is a Mersenne prime.

\term{Perfect Field}
A field $K$ such that every irreducible polynomial over $K$ has distinct roots in its splitting field. All fields of characteristic 0 and all finite fields are perfect.

\term{Perfect Matching}
In graph theory, a set of edges such that every vertex of the graph is incident to exactly one edge in the set. Also called a 1-factor.

\term{Permutation}
A bijection from a set to itself. A permutation of a finite set of $n$ elements is one of the $n!$ possible linear orderings.

\term{Permutation Group}
A group whose elements are permutations of a given set and whose operation is the composition of permutations. The symmetric group $S_n$ is the group of all permutations of $n$ elements.

\term{Permutation Matrix}
A square binary matrix that has exactly one entry of 1 in each row and each column and 0s elsewhere. Multiplying a matrix by a permutation matrix reorders its rows or columns.

\term{Perron--Frobenius Theorem}
A non-negative irreducible square matrix has a positive real eigenvalue equal to its spectral radius, with a strictly positive eigenvector; this eigenvalue is simple and dominates all others in modulus. It is the foundation for the convergence of Markov chains and for stable states in population dynamics, such as those modeled by Leslie matrices.

\term{Phase Space}
In classical mechanics, a space in which all possible states of a system are represented, with plot coordinates asCanadian position and momentum. The volume of a region in phase space is preserved by Hamiltonian flow (Liouville's theorem).

\term{Phi ($\phi$) (Golden Ratio)}
The irrational number $\phi = \frac{1+\sqrt{5}}{2} \approx 1.618$, which is the positive solution to $x^2 - x - 1 = 0$. It appears in the geometry of pentagons, the Fibonacci sequence, and phyllotaxis.

\term{Picard--Lindelöf Theorem}
For the initial value problem $y' = f(x,y)$, $y(x_0) = y_0$, if $f$ is continuous in $x$ and Lipschitz continuous in $y$, then a unique solution exists on some interval about $x_0$. The proof typically proceeds via the contraction mapping principle applied to a space of continuous functions.

\term{Pigeonhole Principle}
A combinatorial principle stating that if $n$ items are put into $m$ containers, with $n > m$, then at least one container must contain more than one item.

\term{Pointwise Convergence}
A sequence of functions $f_n$ converges pointwise to $f$ on a set $S$ if for every $x \in S$, the sequence of values $f_n(x)$ converges to $f(x)$. This is a weaker condition than uniform convergence.

\term{Poisson Distribution}
A discrete probability distribution that expresses the probability of a given number of events occurring in a fixed interval of time or space, provided these events occur with a known constant mean rate and independently of the time since the last event.

\term{Poisson Equation}
A partial differential equation of the form $\Delta u = f$, where $\Delta$ is the Laplacian. It describes the potential $u$ generated by a source density $f$, appearing in electrostatics and gravity.

\term{Poisson Process}
A stochastic process $N(t)$ that counts the number of events occurring up to time $t$, where the increments are independent and the number of events in any interval follows a Poisson distribution.

\term{Poisson Summation Formula}
A relationship between the sum of values of a function at integers and the sum of values of its Fourier transform at integers: $\sum_{n=-\infty}^\infty f(n) = \sum_{k=-\infty}^\infty \hat{f}(k)$.

\term{Polytope}
The generalization of a polygon (2D) and polyhedron (3D) to $n$ dimensions. A convex polytope is the convex hull of a finite set of points in $\mathbb{R}^n$.

\term{Polynomial}
An expression consisting of variables and coefficients, using only addition, subtraction, multiplication, and non-negative integer exponents of variables.

\term{Polynomial Ring}
The ring $R[x]$ consisting of all polynomials with coefficients in a ring $R$. If $R$ is a field, $R[x]$ is a Euclidean domain.

\term{Polar Coordinates}
A two-dimensional coordinate system in which each point is determined by a distance $r$ from a reference point (the origin) and an angle $\theta$ from a reference direction.

\term{Polar Decomposition}
The factorization of a complex square matrix $A$ as $A = UP$, where $U$ is a unitary matrix and $P$ is a positive semi-definite Hermitian matrix.

\term{Pontryagin Duality}
A duality between a locally compact abelian (LCA) group $G$ and its dual group $\widehat{G}$, which consists of all continuous characters of $G$. The double dual is canonically isomorphic to $G$.

\term{Poincaré Disk Model}
A model of hyperbolic geometry where the hyperbolic plane is represented by the interior of the unit disk in $\mathbb{R}^2$. Geodesics are arcs of circles orthogonal to the boundary.

\term{Poincaré Recurrence Theorem}
A theorem in dynamical systems stating that for a measure-preserving transformation on a finite-measure space, almost every point in any set $A$ will eventually return to $A$ infinitely often.

\term{Poincaré Conjecture}
A famous theorem stating that every simply connected, closed 3-manifold is homeomorphic to the 3-sphere. It was proven by Grigori Perelman (2003) using Ricci flow.

\term{Poincaré Duality}
A fundamental result in manifold topology: for a compact oriented $n$-manifold $M$, there is an isomorphism $H^k(M; \mathbb{Z}) \cong H_{n-k}(M; \mathbb{Z})$.

\term{Population}
The entire set of objects of interest in a statistical study, from which samples are drawn.

\term{Power Set}
The set of all subsets of a set $S$, denoted $\mathcal{P}(S)$ or $2^S$. By Cantor's theorem, the cardinality of the power set is strictly greater than the cardinality of $S$.

\term{Power Series}
An infinite series of the form $\sum_{n=0}^\infty a_n (x-c)^n$. Each power series has a radius of convergence $R$, within which the series converges absolutely.

\term{Preconditioner}
A matrix or operator that transforms a linear system into one with better conditioning for iterative solvers.

\term{Preimage}
For a function $f: X \to Y$ and a subset $B \subseteq Y$, the preimage $f^{-1}(B)$ is the set of all $x \in X$ such that $f(x) \in B$.

\term{Pre-order}
A binary relation that is reflexive and transitive, but not necessarily antisymmetric. A partial order is a pre-order that is also antisymmetric.

\term{Presheaf}
An assignment of data to open sets together with restriction maps, compatible with inclusions of open sets.

\term{Primal Problem}
The original optimization problem, to which a dual problem is associated.

\term{Prime Ideal}
An ideal $P$ such that $ab \in P$ implies $a \in P$ or $b \in P$; the quotient ring is an integral domain.

\term{Prime Number Theorem}
The number $\pi(x)$ of primes at most $x$ is asymptotically $\pi(x) \sim x / \ln x$. An equivalent statement is that $\pi(x)$ is asymptotic to the logarithmic integral $\operatorname{li}(x)$, which approximates $\pi(x)$ even more closely. The size of the error term depends on the zeros of the Riemann zeta function, and the Riemann hypothesis would imply a sharp square-root error bound.

\term{Principal Ideal}
An ideal $I$ in a ring $R$ that is generated by a single element $a$: $I = (a) = \{ra : r \in R\}$.

\term{Principal Component Analysis (PCA)}
A statistical procedure that uses an orthogonal transformation to convert a set of observations of possibly correlated variables into a set of values of linearly uncorrelated variables called principal components.

\term{Principal Fiber Bundle}
A fiber bundle $\pi: P \to B$ together with a Lie group $G$ that acts freely and transitively on each fiber. These are used to define gauge theories in physics.

\term{Principal Ideal Domain (PID)}
An integral domain in which every ideal is a principal ideal. Examples include the ring of integers $\mathbb{Z}$ and the ring of polynomials $F[x]$ over a field $F$.

\term{Principal Value (Cauchy)}al
The Cauchy principal value is a method for assigning a value to certain divergent integrals, such as $\int_{-a}^a \frac{1}{x}\,dx$, by taking a symmetric limit around the singularity.

\term{Probability Density Function}
The function $f$ such that $P(a \leq X \leq b) = \int_a^b f(x)\,dx$ for a continuous random variable $X$.

\term{Probability Space}
A mathematical triple $(\Omega, \mathcal{F}, P)$ consisting of a sample space $\Omega$, a $\sigma$-algebra $\mathcal{F}$ of events, and a probability measure $P$.

\term{Product Space}
The Cartesian product of two or more sets, spaces, or groups, equipped with the product structure (e.g., the product topology or product measure).

\term{Product Topology}
The coarsest topology on a product $\prod X_i$ such that all projection maps $\pi_i$ are continuous. A basis for this topology consists of products of open sets, where all but finitely many are the whole space $X_i$.

\term{Projection (Linear Algebra)}
A linear operator $P: V \to V$ such that $P^2 = P$. It decomposes the space $V$ into the direct sum of the image of $P$ and the kernel of $P$.

\term{Projective Geometry}
The study of geometric properties that are invariant under projective transformations. It extends Euclidean geometry by adding "points at infinity."

\term{Projective Space}
The set of all lines passing through the origin in a vector space $\mathbb{R}^{n+1}$. It is denoted $\mathbb{P}^n(\mathbb{R})$ and can be viewed as the quotient of $\mathbb{R}^{n+1} \setminus \{0\}$ by the equivalence $v \sim \lambda v$.

\term{Projective Variety}
The set of common zeros of a collection of homogeneous polynomials in a projective space. Projective varieties are the central objects of study in classical algebraic geometry.

\term{Proper Map}
A continuous map $f: X \to Y$ between topological spaces is proper if the preimage of every compact set in $Y$ is compact in $X$.

\term{Primitive Root}
An integer $g$ is a primitive root modulo $n$ if every integer $a$ coprime to $n$ is congruent to a power of $g$ modulo $n$. This means $g$ generates the multiplicative group $(\mathbb{Z}/n\mathbb{Z})^\times$.

\term{Primitive Element Theorem}
A theorem in field theory stating that every finite separable extension $L/K$ is a simple extension, meaning $L = K(\alpha)$ for some $\alpha \in L$.

\term{Proportionality}
A relationship between two quantities where their ratio remains constant. If $y$ is proportional to $x$, then $y = kx$ for some constant $k$.

\term{Proximal Operator}
The map $\operatorname{prox}_f(x) = \arg\min_y (f(y) + \frac{1}{2}\|x-y\|^2)$, central to proximal methods in optimization.

\term{Pullback}
The map $f^*$ induced by a function $f$, sending objects on the codomain to objects on the domain, such as differential forms or bundles.

\term{Purely Imaginary}
A complex number of the form $bi$ with real part zero.

\term{Pushforward}
The map $f_*$ induced by a function $f$, pushing vectors from the tangent space of the domain to that of the codomain.
